% অনুশীলনী ভান্ডার — পাঠ ২
% পাঠটি এই ফাইলের নাম থেকে নির্ধারিত হয়।

\begin{exo}[ঠান্ডা বরফ থেকে ফুটন্ত জল: মাত্রার ক্রম]{chauffage-glace-eau-ebullition}
\seoready{true}
\keywords{ক্যালরিমিতি, দশা পরিবর্তন, গলন, বাষ্পীভবন, সুপ্ত তাপ, মাত্রার ক্রম}

\textbf{সাংখ্যিক উপাত্ত (বায়ুমণ্ডলীয় চাপে):}
\[
\begin{aligned}
c_{\mathrm{i}} &= 2{,}1\times10^3\,\text{J}\!\cdot\!\text{kg}^{-1}\!\cdot\!\text{K}^{-1},
& L_f &= 3{,}34\times10^5\,\text{J}\!\cdot\!\text{kg}^{-1},\\
c_{\mathrm{w}} &= 4{,}18\times10^3\,\text{J}\!\cdot\!\text{kg}^{-1}\!\cdot\!\text{K}^{-1},
& L_v &= 2{,}26\times10^6\,\text{J}\!\cdot\!\text{kg}^{-1}.
\end{aligned}
\]
বায়ুমণ্ডলীয় চাপে, প্রাথমিকভাবে $-100\,{}^\circ\text{C}$ তাপমাত্রায় থাকা এক কিলোগ্রাম বরফকে ধীরে ধীরে গরম করা হয়।

\begin{enumerate}
  \item বরফকে $-100\,{}^\circ\text{C}$ থেকে $0\,{}^\circ\text{C}$ পর্যন্ত আনতে প্রয়োজনীয় শক্তি নির্ণয় করুন।
  \item $0\,{}^\circ\text{C}$ তাপমাত্রায় বরফ সম্পূর্ণ গলাতে প্রয়োজনীয় শক্তি নির্ণয় করুন।
  \item তরল জলকে $0\,{}^\circ\text{C}$ থেকে $100\,{}^\circ\text{C}$ পর্যন্ত গরম করতে প্রয়োজনীয় শক্তি নির্ণয় করুন।
  \item $100\,{}^\circ\text{C}$ তাপমাত্রায় জল সম্পূর্ণ বাষ্পীভূত করতে অতিরিক্ত কত শক্তি প্রয়োজন তা নির্ণয় করুন।
  \item কোন ধাপে সর্বাধিক শক্তি লাগে? এক কিলোগ্রাম জলকে $0\,{}^\circ\text{C}$ থেকে $100\,{}^\circ\text{C}$ পর্যন্ত গরম করার শক্তির সঙ্গে দশ মিটার উচ্চতা থেকে $m$ ভরের পতনে মুক্ত বিভব শক্তির তুলনা করুন। একই শক্তি পেতে কত ভর $m$ ফেলতে হবে?
  \item এবার বিপরীত প্রক্রিয়াটি বিবেচনা করুন। $m_v$ ভরের জলীয় বাষ্প একটি গাড়ির সামনের কাচে ঘনীভূত হয়; পরে উৎপন্ন তরল জল আর ঠান্ডা হয় না। ঘনীভবনে মুক্ত তাপ $Q_{\mathrm{rel}}$-এর রাশি লিখুন এবং $m_v=1{,}0\times10^{-2}\,\text{kg}$-এর জন্য তা নির্ণয় করুন। কাচের তাপমাত্রায় $L_v=2{,}45\times10^6\,\text{J}\!\cdot\!\text{kg}^{-1}$ নিন।
\end{enumerate}

\begin{indication}
দশা পরিবর্তন ছাড়া কোনো দশাকে গরম করতে $Q=mc\Delta T$ এবং স্থির তাপমাত্রায় দশা পরিবর্তনের জন্য $Q=mL$ ব্যবহার করুন। $h$ উচ্চতায় $m$ ভরের মহাকর্ষীয় বিভব শক্তি $E_p=mgh$; $g=9{,}81\,\text{m}\!\cdot\!\text{s}^{-2}$ নিন।
\end{indication}

\begin{solution}
\textbf{1.} বরফকে গলনাঙ্ক পর্যন্ত গরম করতে লাগে
\[
Q_1=mc_{\mathrm{i}}\Delta T
=1{,}0\times2{,}1\times10^3\times100
=2{,}1\times10^5\,\text{J}.
\]

\textbf{2.} গলনের সময় তাপমাত্রা $0\,{}^\circ\text{C}$ থাকে:
\[
Q_2=mL_f=1{,}0\times3{,}34\times10^5=3{,}34\times10^5\,\text{J}.
\]

\textbf{3.} তরল জলকে $0$ থেকে $100\,{}^\circ\text{C}$ পর্যন্ত গরম করতে
\[
Q_3=mc_{\mathrm{w}}\Delta T
=1{,}0\times4{,}18\times10^3\times100
=4{,}18\times10^5\,\text{J}.
\]

\textbf{4.} সম্পূর্ণ বাষ্পীভবনে আরও লাগে
\[
Q_4=mL_v=1{,}0\times2{,}26\times10^6=2{,}26\times10^6\,\text{J}.
\]
অতএব শুধু এই ধাপেই কয়েক মেগাজুল মাত্রার শক্তি লাগে।

\textbf{5.} বাষ্পীভবনই নিঃসন্দেহে সবচেয়ে শক্তিব্যয়ী ধাপ:
\[
Q_4=2{,}26\times10^6\,\text{J}
>Q_3=4{,}18\times10^5\,\text{J}
>Q_2=3{,}34\times10^5\,\text{J}
>Q_1=2{,}1\times10^5\,\text{J}.
\]
বিশেষত, তরল জলকে $0$ থেকে $100\,{}^\circ\text{C}$ পর্যন্ত গরম করার তুলনায় বাষ্পীভবনে প্রায় $(2{,}26\times10^6)/(4{,}18\times10^5)\approx5{,}4$ গুণ বেশি শক্তি লাগে।

$h=10\,\text{m}$ উচ্চতা থেকে $m$ ভর পড়লে মুক্ত বিভব শক্তি $E_p=mgh$। $mgh=Q_3$ বসিয়ে পাই
\[
m=\frac{Q_3}{gh}
=\frac{4{,}18\times10^5}{9{,}81\times10}
\approx4{,}26\times10^3\,\text{kg}.
\]
অতএব এক কিলোগ্রাম জলকে $0$ থেকে $100\,{}^\circ\text{C}$ পর্যন্ত গরম করার সমান শক্তি মুক্ত করতে দশ মিটার উচ্চতা থেকে প্রায় $4{,}3$ টন ভর ফেলতে হবে! এটি বিপুল; তাই গৃহস্থালির শক্তি ব্যবহারে জল গরম করার অংশটি গুরুত্বপূর্ণ। জলের ভর-নির্দিষ্ট তাপধারণ ক্ষমতাও প্রচলিত ধাতুগুলির তুলনায় খুব বেশি (যেমন তামার প্রায় দশ গুণ; পরের অনুশীলনগুলি দেখুন)।

\textbf{6.} ঘনীভবন বাষ্পীভবনের বিপরীত প্রক্রিয়া: বাষ্প কাচে সুপ্ত তাপ ছেড়ে দেয়,
\[
Q_{\mathrm{rel}}=m_vL_v.
\]
$m_v=1{,}0\times10^{-2}\,\text{kg}$ হলে
\[
Q_{\mathrm{rel}}=1{,}0\times10^{-2}\times2{,}45\times10^6
=2{,}45\times10^4\,\text{J}.
\]
অতএব অল্প ভরের বাষ্পের ঘনীভবনও উপেক্ষণীয় নয় এমন তাপ মুক্ত করতে পারে। এই তাপ কাচ ও তার পরিবেশ গ্রহণ করে।
\end{solution}
\end{exo}

\begin{exo}[একটি বড় গাছের বাষ্পোৎসর্জন: প্রাকৃতিক শীতাতপনিয়ন্ত্রক?]{odg-evapotranspiration-arbre}
\seoready{true}
\keywords{মাত্রার ক্রম, বাষ্পোৎসর্জন, গাছ, সুপ্ত তাপ, শীতলীকরণ ক্ষমতা, শীতাতপনিয়ন্ত্রক, COP}

গরম ও শুষ্ক দিনে পর্যাপ্ত জল পাওয়া একটি বড় গাছ প্রায় $500\,\text{L}=5{,}0\times10^{-1}\,\text{m}^3$ জল বাষ্পোৎসর্জন করতে পারে। বাষ্পোৎসর্জন প্রধানত আলোতে ঘটে বলে ধরে নিন যে তা বারো ঘণ্টা জুড়ে বিস্তৃত। গাছের পাতার ছাউনি ভূমিতে $A_{\mathrm{tree}}=50\,\text{m}^2$ ক্ষেত্রফল ঢেকে রাখে। বায়ুমণ্ডলীয় চাপে জলের ঘনত্ব ও বাষ্পীভবনের সুপ্ত তাপ:
\[
\rho_{\mathrm{w}}=1{,}0\times10^3\,\text{kg}\!\cdot\!\text{m}^{-3},
\qquad L_v=2{,}45\times10^6\,\text{J}\!\cdot\!\text{kg}^{-1}.
\]

\begin{enumerate}
  \item দিনে গাছটি যত জল বাষ্পোৎসর্জন করে তার ভর এবং সেই জল বাষ্পীভূত করতে প্রয়োজনীয় শক্তি নির্ণয় করুন।
  \item এতে শীতলীকরণ ঘটে কেন ব্যাখ্যা করুন।
  \item সক্রিয় বাষ্পোৎসর্জনের বারো ঘণ্টায় প্রতি বর্গমিটারে গাছটির গড় শীতলীকরণ ক্ষমতা নির্ণয় করুন।
  \item তুলনার জন্য $P_{\mathrm{elec}}=2{,}0\times10^3\,\text{W}$ বৈদ্যুতিক ক্ষমতায় চালিত, কার্যক্ষমতা গুণাঙ্ক $\mathrm{COP}=3{,}0$-এর একটি শীতাতপনিয়ন্ত্রক $A_{\mathrm{room}}=20\,\text{m}^2$ ক্ষেত্রফলের একটি কক্ষ ঠান্ডা করে। সংজ্ঞা অনুসারে
\[
\mathrm{COP}=\frac{P_{\mathrm{cool}}}{P_{\mathrm{elec}}},
\]
যেখানে $P_{\mathrm{cool}}$ কক্ষ থেকে অপসারিত তাপীয় ক্ষমতা। শীতাতপনিয়ন্ত্রকের $P_{\mathrm{cool}}$ এবং প্রতি বর্গমিটারে শীতলীকরণ ক্ষমতা নির্ণয় করে গাছের ফলের সঙ্গে তুলনা করুন।
  \item একক ক্ষেত্রফলে ক্ষমতার মাত্রার ক্রম একই হলেও গাছ ও শীতাতপনিয়ন্ত্রক ঠিক একই অনুভূত শীতলতা দেয়—এ সিদ্ধান্ত কেন নেওয়া যায় না?
  \item আপনার মতে অনেক ঘরোয়া উদ্ভিদ দিয়ে কি বাড়ি ঠান্ডা করা যায়? (সাধারণ জ্ঞান: উত্তরটি জৈবিক প্রক্রিয়ার উপর নির্ভরশীল।)
\end{enumerate}

\begin{indication}
$m=\rho V$, $Q_{\mathrm{evap}}=mL_v$ এবং $P=Q/\tau$ ব্যবহার করুন। শীতাতপনিয়ন্ত্রকের জন্য $P_{\mathrm{cool}}=\mathrm{COP}\,P_{\mathrm{elec}}$।
\end{indication}

\begin{solution}
\textbf{1.} $500\,\text{L}=5{,}0\times10^{-1}\,\text{m}^3$, তাই
\[
m=\rho_{\mathrm{w}}V=1{,}0\times10^3\times5{,}0\times10^{-1}=5{,}0\times10^2\,\text{kg},
\]
এবং
\[
Q_{\mathrm{evap}}=mL_v=5{,}0\times10^2\times2{,}45\times10^6=1{,}225\times10^9\,\text{J}.
\]
এর মাত্রার ক্রম এক গিগাজুল।

\textbf{2.} বাষ্পীভবন একটি তাপগ্রাহী দশা পরিবর্তন। এর জন্য প্রয়োজনীয় $Q_{\mathrm{evap}}$ পাতা, মাটি ও চারপাশের বায়ু থেকে নেওয়া হয়; ফলে তাদের তাপমাত্রা কমে এবং শীতলীকরণ ঘটে।

\textbf{3.} বারো ঘণ্টা মানে $\tau=12\times3600=4{,}32\times10^4\,\text{s}$। অতএব
\[
P_{\mathrm{tree}}=\frac{Q_{\mathrm{evap}}}{\tau}
=\frac{1{,}225\times10^9}{4{,}32\times10^4}
\approx2{,}84\times10^4\,\text{W},
\]
এবং পাতার ছাউনির $50\,\text{m}^2$ ক্ষেত্রফলে
\[
\frac{P_{\mathrm{tree}}}{A_{\mathrm{tree}}}\approx5{,}7\times10^2\,\text{W}\!\cdot\!\text{m}^{-2}.
\]

\textbf{4.} উপযোগী শীতলীকরণ ক্ষমতা
\[
P_{\mathrm{cool}}=\mathrm{COP}\,P_{\mathrm{elec}}=3{,}0\times2{,}0\times10^3=6{,}0\times10^3\,\text{W}.
\]
কক্ষের ক্ষেত্রফল অনুসারে
\[
\frac{P_{\mathrm{cool}}}{A_{\mathrm{room}}}=\frac{6{,}0\times10^3}{20}=3{,}0\times10^2\,\text{W}\!\cdot\!\text{m}^{-2}.
\]
এই মডেলে বাষ্পোৎসর্জনের ক্ষমতা একটি সাধারণ শীতাতপনিয়ন্ত্রকের প্রায় দ্বিগুণ।

\textbf{5.} তুলনাটি শুধু শক্তির প্রবাহ নিয়ে। শীতাতপনিয়ন্ত্রক একটি বদ্ধ, নিয়ন্ত্রিত আয়তন ঠান্ডা করে; গাছের বাষ্প বায়ুমণ্ডলে ছড়িয়ে যায় এবং বাতাস আবার গরম বায়ু আনে। অনুভূত প্রভাব বায়ু চলাচল, আর্দ্রতা, পরিবেশের তাপমাত্রা ও গাছ থেকে দূরত্বের উপর নির্ভর করে। সূর্যালোক, বাতাস, আর্দ্রতা, প্রজাতি ও জলের প্রাপ্যতার সঙ্গে বাষ্পোৎসর্জনও বদলায় এবং খরায় খুব কমে যেতে পারে। অন্যদিকে গাছ ছায়াও দেয়। তাই হিসাবগুলি মাত্রার ক্রমের তুলনা দেয়, সমান তাপীয় আরামের নিশ্চয়তা নয়।

\textbf{6.} বাস্তবে ঘরোয়া উদ্ভিদের শীতলীকরণ খুব সামান্য হবে। অধিকাংশ উদ্ভিদের আলোতে পত্ররন্ধ্র খোলে এবং জলীয় বাষ্প বের হয়; ঘরের কম আলোতে পত্ররন্ধ্র কম খোলে। কম বায়ু চলাচলের ঘরে আর্দ্রতা বাড়লেও বাষ্পীভবন ধীরে যায়। অনেক উদ্ভিদ অল্প স্থানীয় বাষ্পীভবনজনিত শীতলতা দিলেও শীতাতপনিয়ন্ত্রকের বিকল্প নয়। আর্দ্র বায়ু বাইরে পাঠাতে বায়ু নবায়ন দরকার। তীব্র কৃত্রিম আলো বাষ্পোৎসর্জন বাড়াতে পারে, কিন্তু তার বৈদ্যুতিক শক্তি শেষ পর্যন্ত ঘরেই প্রায় পুরোটা তাপে পরিণত হয়। ক্যাকটাসসহ শুষ্ক পরিবেশে অভিযোজিত CAM উদ্ভিদ ব্যতিক্রম: তাদের পত্ররন্ধ্র প্রধানত রাতে খোলে।
\end{solution}
\end{exo}

\begin{exo}[দুই ভরের জলের মিশ্রণ]{calorimetrie-melange-eau}
\seoready{true}
\keywords{ক্যালরিমিতি, মিশ্রণ, তাপধারণ ক্ষমতা, বিশিষ্ট তাপ, জোসেফ ব্ল্যাক, সাম্য তাপমাত্রা}

$T_1=80\,{}^\circ\text{C}$ তাপমাত্রার $m_1=0{,}200\,\text{kg}$ জল এমন একটি পাত্রে ঢালা হলো যাতে আগে থেকেই $T_2=15\,{}^\circ\text{C}$ তাপমাত্রার $m_2=0{,}300\,\text{kg}$ জল আছে। পাত্রটি আদর্শ ক্যালরিমিটার: বাইরের দিকে তাপক্ষয় এবং পাত্রের নিজস্ব তাপধারণ ক্ষমতা উপেক্ষা করা হয়। সম্পর্কটি হলো
\[
Q=mc\Delta T,
\]
যেখানে তরল জলের ভর-নির্দিষ্ট তাপধারণ ক্ষমতা $c=4{,}18\times10^3\,\text{J}\!\cdot\!\text{kg}^{-1}\!\cdot\!\text{K}^{-1}$।

\begin{enumerate}
  \item গরম জল যে তাপ ছাড়ে ঠান্ডা জল তা সম্পূর্ণ গ্রহণ করে—এটি ব্যাখ্যা করে $m_1$, $m_2$, $c$, $T_1$, $T_2$ ও সাম্য তাপমাত্রা $T_{eq}$-সহ সংরক্ষণ সমীকরণ লিখুন।
  \item $T_{eq}$-এর বীজগাণিতিক রাশি ও সাংখ্যিক মান নির্ণয় করুন।
  \item মিশ্রণের সময় গরম জল যে তাপ $Q$ ছাড়ে তা নির্ণয় করুন।
  \item একই পদার্থের এমন মিশ্রণ দিয়ে কি $c$ নির্ণয় করা যায়? যুক্তি দিন।
\end{enumerate}

\begin{indication}
ক্যালরিমিটারটি অন্তরিত ও অনমনীয় হওয়ায় মোট অভ্যন্তরীণ শক্তি $U_1+U_2$ সংরক্ষিত: একটি অংশ যে তাপ ছাড়ে, অন্যটি চিহ্নের বিপরীতে সেই তাপ গ্রহণ করে।
\end{indication}

\begin{solution}
\textbf{1.} অন্তরিত ক্যালরিমিটারে গরম জল ও ঠান্ডা জলের তন্ত্রের মোট অভ্যন্তরীণ শক্তি সংরক্ষিত। প্রতিটি জলভরের প্রাথমিক অবস্থা থেকে সাম্য পর্যন্ত $Q=mc\Delta T$ প্রয়োগ করলে
\[
m_1c(T_{eq}-T_1)+m_2c(T_{eq}-T_2)=0.
\]

\textbf{2.} একই পদার্থ হওয়ায় $c$ কেটে যায়:
\[
T_{eq}=\frac{m_1T_1+m_2T_2}{m_1+m_2}
=\frac{0{,}200\times80+0{,}300\times15}{0{,}200+0{,}300}
=41\,{}^\circ\text{C}.
\]

\textbf{3.} গরম জল যে তাপ ছাড়ে
\[
Q=m_1c(T_1-T_{eq})
=0{,}200\times4{,}18\times10^3\times(80-41)
\approx3{,}26\times10^4\,\text{J}.
\]
ঠান্ডা জলের গ্রহণ করা তাপও একই:
\[
m_2c(T_{eq}-T_2)=0{,}300\times4{,}18\times10^3\times26\approx3{,}26\times10^4\,\text{J}.
\]

\textbf{4.} না। একই পদার্থের দুই ভরের জন্য শক্তির হিসাবে একই $c$ দুই পদেই গুণক হয়ে কেটে যায়। তাই মাপা সাম্য তাপমাত্রা $c$-এর উপর নির্ভর করে না; এটি শুধু ভর-দ্বারা ভারিত প্রাথমিক তাপমাত্রার গড়। মিশ্রণ পদ্ধতিতে ভর-নির্দিষ্ট তাপধারণ ক্ষমতা মাপতে পরের অনুশীলনের মতো পরিচিত তাপধারণ ক্ষমতার একটি পদার্থ দরকার।
\end{solution}
\end{exo}

\begin{exo}[মিশ্রণ পদ্ধতিতে একটি ধাতুর ভর-নির্দিষ্ট তাপধারণ ক্ষমতা নির্ণয়]{calorimetrie-methode-melanges}
\seoready{true}
\keywords{ক্যালরিমিতি, মিশ্রণ পদ্ধতি, ভর-নির্দিষ্ট তাপধারণ ক্ষমতা, ক্যালরিমিটার, বিশিষ্ট তাপ, জল-সমতুল্য}

XVIII\textsuperscript{e} শতাব্দীতে জোসেফ ব্ল্যাক যেমন করতেন, তেমন \emph{মিশ্রণ পদ্ধতিতে} অজানা ধাতুর ভর-নির্দিষ্ট তাপধারণ ক্ষমতা $c_m$ মাপতে চাই। $m=0{,}150\,\text{kg}$ নমুনাটিকে $T_m=95\,{}^\circ\text{C}$ পর্যন্ত গরম করে দ্রুত এমন ক্যালরিমিটারে ডোবানো হয় যাতে $T_e=18\,{}^\circ\text{C}$ তাপমাত্রার $m_e=0{,}250\,\text{kg}$ জল আছে; $c_e=4{,}18\times10^3\,\text{J}\!\cdot\!\text{kg}^{-1}\!\cdot\!\text{K}^{-1}$। তাপীয় সাম্যে $T_{eq}=22\,{}^\circ\text{C}$ মাপা হয়। প্রথমে ক্যালরিমিটার পাত্রের (দেয়াল, নাড়নি, থার্মোমিটার) নিজস্ব তাপধারণ ক্ষমতা উপেক্ষা করুন।

\begin{enumerate}
  \item আগের অনুশীলনের মতো ধাতু ও জলের অন্তরিত তন্ত্রের শক্তি সংরক্ষণ লিখুন, যাতে $c_m$ একমাত্র অজানা থাকে।
  \item $c_m$-এর বীজগাণিতিক রাশি ও সাংখ্যিক মান নির্ণয় করুন। এটি আনুমানিক কোন পরিচিত ধাতুর সঙ্গে মেলে? তুলনার মান: অ্যালুমিনিয়াম $\approx9{,}0\times10^2\,\text{J}\!\cdot\!\text{kg}^{-1}\!\cdot\!\text{K}^{-1}$, লোহা $\approx4{,}5\times10^2$, তামা $\approx3{,}9\times10^2$, সিসা $\approx1{,}3\times10^2$।
  \item বাস্তবে ক্যালরিমিটার পাত্রও ধাতুর ছাড়া তাপের একটি অংশ গ্রহণ করে। একে \emph{জল-সমতুল্য} $\mu=1{,}5\times10^{-2}\,\text{kg}$ দিয়ে মডেল করা হয়: ক্যালরিমিটারটি যেন অতিরিক্ত $\mu$ ভরের জলের মতো আচরণ করে। $\mu$ ধরে $c_m$ আবার নির্ণয় করুন এবং সংশোধনের দিক ব্যাখ্যা করুন (সংশোধিত মান প্রশ্ন ২-এর চেয়ে বড় না ছোট?)।
  \item পরীক্ষায় নমুনাটি \emph{দ্রুত} ক্যালরিমিটারে ডোবানো এবং মাপার সময় পরিবেশের বায়ু থেকে ভালোভাবে তাপীয়ভাবে অন্তরিত রাখা জরুরি কেন?
\end{enumerate}

\begin{indication}
ধাতু $Q_m=mc_m(T_m-T_{eq})$ তাপ ছাড়ে; তা জল (এবং প্রশ্ন ৩-এ অতিরিক্ত $\mu$ জলভর-সমতুল্য ক্যালরিমিটার) সম্পূর্ণ গ্রহণ করে: $Q_m=(m_e+\mu)c_e(T_{eq}-T_e)$।
\end{indication}

\begin{solution}
\textbf{1.} অন্তরিত তন্ত্রে ধাতু $T_m$ থেকে $T_{eq}$-এ ঠান্ডা হয়ে যে তাপ ছাড়ে, জল $T_e$ থেকে $T_{eq}$-এ গরম হয়ে তা সম্পূর্ণ গ্রহণ করে:
\[
mc_m(T_m-T_{eq})=m_ec_e(T_{eq}-T_e).
\]

\textbf{2.}
\[
c_m=\frac{m_ec_e(T_{eq}-T_e)}{m(T_m-T_{eq})}
=\frac{0{,}250\times4{,}18\times10^3\times(22-18)}{0{,}150\times(95-22)}
\approx3{,}82\times10^2\,\text{J}\!\cdot\!\text{kg}^{-1}\!\cdot\!\text{K}^{-1}.
\]
এটি তামার মানের ($\approx3{,}9\times10^2\,\text{J}\!\cdot\!\text{kg}^{-1}\!\cdot\!\text{K}^{-1}$) খুব কাছাকাছি।

\textbf{3.} জল ও পাত্র একসঙ্গে তাপ গ্রহণ করে:
\[
c_m=\frac{(m_e+\mu)c_e(T_{eq}-T_e)}{m(T_m-T_{eq})}
=\frac{(0{,}250+0{,}015)\times4{,}18\times10^3\times4}{0{,}150\times73}
\approx4{,}05\times10^2\,\text{J}\!\cdot\!\text{kg}^{-1}\!\cdot\!\text{K}^{-1}.
\]
সংশোধনে $c_m$ সামান্য বাড়ে। প্রশ্ন ২-এ পাত্রের জল-সমতুল্য উপেক্ষা করায় ধাতুর সত্যিকারের ছাড়া তাপ—যার একটি অংশ পাত্র গরম করেছিল—এবং তাই $c_m$ কম ধরা হয়েছিল।

\textbf{4.} পদ্ধতিটি মাপার সময় তন্ত্রটি অন্তরিত—এই অনুমানের উপর দাঁড়ায়। দ্রুত ডোবালে জলে পৌঁছানোর আগে নমুনার সঙ্গে বায়ুর তাপ বিনিময়ের সময় কমে; ভালো অন্তরণ সাম্য প্রতিষ্ঠার সময় বাইরের তাপক্ষয় কমায়। হিসাবের বাইরে যেকোনো তাপ বিনিময় প্রশ্ন ১-এর শক্তি হিসাব এবং সেখান থেকে পাওয়া $c_m$ ভুল করবে। ব্ল্যাক এবং পরে লাভয়জিয়ে ও লাপ্লাস তাঁদের বরফ-ক্যালরিমিটারে ঠিক এমন পরীক্ষামূলক সতর্কতাই নিয়েছিলেন।
\end{solution}
\end{exo}

\begin{exo}[খাদ্য ক্যালরি ও বিপাকীয় ক্ষমতা]{calorimetrie-calories-alimentaires}
\seoready{true}
\keywords{ক্যালরি, কিলোক্যালরি, খাদ্য ক্যালরি, যান্ত্রিক সমতুল্য, বিপাকীয় ক্ষমতা, একক}

পুষ্টি-তথ্যে বলা হয়, একজন প্রাপ্তবয়স্ক প্রতিদিন গড়ে প্রায় $2000\,\text{kcal}$ গ্রহণ করেন। কিলোক্যালরি পুষ্টিবিদ্যায় এখনও ব্যবহৃত একটি অ-SI একক; এর SI রূপান্তর $1\,\text{kcal}=4{,}18\times10^3\,\text{J}$।

\begin{enumerate}
  \item দৈনিক শক্তির পরিমাণটি জুলে রূপান্তর করুন।
  \item শক্তিটি ২৪ ঘণ্টায় সমভাবে ব্যবহৃত হয় ধরে প্রাপ্তবয়স্কটির গড় বিপাকীয় ক্ষমতা ওয়াটে নির্ণয় করুন।
  \item একটি এনার্জি বারে ``$210\,\text{kcal}$'' লেখা আছে। একই শক্তি সম্পূর্ণ তাপে রূপান্তরিত হলে, $20\,{}^\circ\text{C}$ তাপমাত্রার কত ভর জলকে স্ফুটনাঙ্কে ($100\,{}^\circ\text{C}$) আনা যাবে? $c_{\mathrm{w}}=4{,}18\times10^3\,\text{J}\!\cdot\!\text{kg}^{-1}\!\cdot\!\text{K}^{-1}$ নিন।
  \item মানবদেহে এই শক্তি বাস্তবে কোন কোন রূপে ব্যবহৃত হয় বলে আপনি মনে করেন? (গুণগত উত্তর)
\end{enumerate}

\begin{indication}
প্রশ্ন ২-এর জন্য $\mathcal P=E/\Delta t$ ব্যবহার করুন। প্রশ্ন ৩-এ প্রথমে লেবেলের শক্তিকে জুলে বদলে $Q=mc\Delta T$ থেকে $m$ আলাদা করুন।
\end{indication}

\begin{solution}
\textbf{1.} SI এককে
\[
E=2000\times4{,}18\times10^3=8{,}36\times10^6\,\text{J}.
\]

\textbf{2.}
\[
\mathcal P=\frac{8{,}36\times10^6}{8{,}64\times10^4}\approx97\,\text{W}.
\]
সারা দিনে একজন মানুষ গড়ে প্রায় একটি জ্বলন্ত ফিলামেন্টের বাল্বের সমান ক্ষমতা ব্যবহার করেন। এই বিস্ময়কর মাত্রার ক্রমটি খাদ্য ক্যালরিকে প্রচলিত ভৌত এককে রূপান্তরের উপযোগিতা দেখায়।

\textbf{3.} বারের শক্তি
\[
Q=210\times4{,}18\times10^3\approx8{,}78\times10^5\,\text{J}.
\]
$\Delta T=80\,\text{K}$ হলে
\[
m=\frac{Q}{c\Delta T}
=\frac{8{,}78\times10^5}{4{,}18\times10^3\times80}
\approx2{,}63\,\text{kg}.
\]
অতএব মাত্রার ক্রমে একটি এনার্জি বার প্রায় $2{,}5\,\text{kg}$ কলের জল ফুটিয়ে তোলার শক্তির সমান।

\textbf{4.} দেহ শক্তি অদৃশ্য করে ``খরচ'' করে না; খাদ্যের রাসায়নিক শক্তিকে রূপান্তর করে। এক অংশ পেশির যান্ত্রিক কাজ, অঙ্গের কার্যকলাপ, কোষের বৈদ্যুতিক নতি বজায় রাখা, অণু পরিবহন এবং কলা নবায়নের রাসায়নিক সংশ্লেষে যায়। আরেক অংশ গ্লাইকোজেন ও চর্বিতে রাসায়নিকভাবে সঞ্চিত হতে পারে। ব্যবহৃত শক্তির অধিকাংশ শেষে তাপ হিসেবে অপচয় হয়, যা পরিবেশে যাওয়ার আগে দেহের তাপমাত্রা বজায় রাখতে সাহায্য করে। সামান্য অংশ বর্জ্যে অবশিষ্ট রাসায়নিক শক্তি হিসেবেও দেহ ছাড়ে।
\end{solution}
\end{exo}
